\section{ПРАКТИЧЕСКАЯ РЕАЛИЗАЦИЯ И РАЗРАБОТКА TELEGRAM-БОТА}

В рамках вариативной части задания по проектной практике был разработан интерактивный программный продукт — Telegram-бот «ЭкоМедиа». Основная задача бота заключается в геймификации процесса экологического просвещения и предоставлении персонализированных рекомендаций по улучшению качества жизни.

\subsection{Выбор языка программирования и библиотек}
Для реализации бота был выбран язык программирования \textbf{Python 3}. Python является де-факто стандартом в современной разработке скриптов, ботов и систем анализа данных благодаря своему лаконичному синтаксису, обширной стандартной библиотеке и огромному сообществу разработчиков.

Взаимодействие с серверами Telegram осуществлялось посредством открытого API (Application Programming Interface). Для работы с Telegram Bot API на уровне языка Python была использована внешняя библиотека \texttt{pyTelegramBotAPI} (также известная как \texttt{telebot}). Данная библиотека реализует синхронный подход к обработке сообщений, что идеально подходит для создания легковесных учебных проектов. Она берет на себя рутину по формированию HTTP-запросов (long polling) и сериализации/десериализации JSON-объектов, позволяя разработчику сосредоточиться исключительно на бизнес-логике приложения.

Для развертывания бота был создан конфигурационный файл \texttt{requirements.txt}, содержащий точные версии зависимостей, что гарантирует воспроизводимость окружения на любом хосте.

\subsection{Проектирование архитектуры и хранение состояний}
Одной из главных сложностей при разработке диалоговых интерфейсов (чат-ботов) является необходимость отслеживать контекст общения с конкретным пользователем. Telegram Bot API является stateless-системой (системой без состояния), то есть каждый входящий запрос рассматривается изолированно от предыдущих.

Для реализации сложных сценариев (например, пошагового прохождения теста на цифровую гигиену) потребовалась реализация паттерна «Конечный автомат» (Finite State Machine, FSM). В архитектуре бота было выделено глобальное хранилище состояний в оперативной памяти в виде Python-словаря (dictionary), ключом в котором выступает уникальный идентификатор чата (\texttt{chat\_id}).

Пользователь в любой момент времени находится в одном из состояний (например, \texttt{STATE\_IDLE}, \texttt{STATE\_QUIZ\_1}, \texttt{STATE\_HYGIENE\_SCREEN\_TIME}). При поступлении нового сообщения бот сначала проверяет текущее состояние пользователя в словаре, а затем направляет сообщение в соответствующую функцию-обработчик (handler). 

\subsection{Реализация алгоритмов и функциональных модулей}
Функционал бота логически разделен на три основных модуля: эко-викторину, тест цифровой гигиены и генератор советов.

\textbf{Модуль 1. Эко-викторина (Eco-Quiz)}\\
Викторина направлена на проверку базовых знаний в области экологии человека (качество воды, микропластик). Модуль содержит заранее заданный массив вопросов. Для взаимодействия с пользователем используется механизм \texttt{InlineKeyboardMarkup} — клавиатура, прикрепленная непосредственно к сообщению.
Когда пользователь нажимает кнопку с вариантом ответа, Telegram отправляет серверу объекта \texttt{CallbackQuery}. Бот перехватывает этот запрос, проверяет правильность ответа, отправляет пользователю уведомление (\texttt{answerCallbackQuery}) и переводит автомат в следующее состояние для отправки нового вопроса. По завершении викторины выводится итоговый счет.

\textbf{Модуль 2. Тест цифровой гигиены}\\
Этот модуль представляет собой систему скоринга (оценки). Бот последовательно задает пользователю вопросы, касающиеся его повседневных привычек:
\begin{enumerate}
    \item Сколько часов в день вы проводите за экранами смартфонов/компьютеров?
    \item Используете ли вы фильтры синего света или ночной режим в вечернее время?
    \item За какое время до сна вы прекращаете использование гаджетов?
\end{enumerate}
Пользователь отвечает с помощью обычных текстовых кнопок (\texttt{ReplyKeyboardMarkup}). В зависимости от выбранного варианта бот начисляет или снимает определенное количество баллов. После прохождения всех этапов алгоритм суммирует баллы и на основе полученного результата (низкий, средний, высокий уровень риска) выдает персонализированное текстовое заключение о состоянии цифровой гигиены пользователя и медицинские рекомендации.

\textbf{Модуль 3. Генератор экологических советов}\\
Для поддержания вовлеченности аудитории реализован модуль выдачи случайных полезных советов. Бза знаний содержит десятки рекомендаций по сортировке мусора, выбору безопасной посуды, минимизации использования пластика и эргономике рабочего места. При вызове команды \texttt{/advice} бот использует модуль \texttt{random} стандартной библиотеки Python для выбора случайного элемента из списка и отправляет его пользователю. Это стимулирует пользователей возвращаться к боту ежедневно.

\subsection{Интеграция с веб-ресурсом}
Одной из задач проекта было объединение всех разработанных компонентов в единую экосистему. Для этого в Telegram-боте реализована команда \texttt{/site}, которая генерирует приветственное сообщение с ссылками на ключевые разделы веб-сайта проекта «ЭкоМедиа». В свою очередь, на сайте размещена специальная страница (\texttt{bot.html}) с описанием функционала бота, QR-кодом и deep-link ссылкой вида \texttt{t.me/ecomediaprobot}, которая позволяет мгновенно открыть диалог с ботом на любом устройстве. Таким образом был реализован принцип кросс-платформенного медиа-потребления.
